\section{Matrix Chain Multiplication}
\label{sec:dp-mcm}

\problemstatement{Given an array $\textit{arr}$ of $n+1$ integers
representing $n$ matrices, where the $i$-th matrix has dimensions
$\textit{arr}[i-1] \times \textit{arr}[i]$, find the
\textbf{minimum number of scalar multiplications} needed to compute
the product of all $n$ matrices (matrix multiplication is associative,
so the order of multiplying \emph{pairs} can be chosen freely).}

\begin{patternbox}
\emph{``Minimum cost to fully combine a chain, where combining two
adjacent groups costs something dependent on their boundaries''} is
this chapter's foundational interval DP: for matrices $i$ through $j$,
try every possible point $k$ to split the chain into a left group and
a right group, multiply those two groups' results together (this
costs $\textit{arr}[i-1] \times \textit{arr}[k] \times \textit{arr}[j]$
scalar multiplications, since the left group reduces to one
$\textit{arr}[i-1]\times\textit{arr}[k]$ matrix and the right to one
$\textit{arr}[k]\times\textit{arr}[j]$ matrix), and take the best
split.
\end{patternbox}

\subsection*{Deriving the recurrence}

Let $f(i,j)$ be the minimum scalar multiplications to multiply matrices
$i$ through $j$ (using $1$-indexed matrix numbers, where matrix $i$ has
dimensions $\textit{arr}[i-1] \times \textit{arr}[i]$). For a single
matrix, $f(i,i)=0$ (nothing to multiply). For $i<j$, try every split
point $k$ from $i$ to $j-1$ (matrices $i..k$ form the left group,
$k{+}1..j$ the right group):
\[
  f(i,j) = \min_{i \le k < j} \Big( f(i,k) + f(k{+}1,j) +
  \textit{arr}[i{-}1] \cdot \textit{arr}[k] \cdot \textit{arr}[j] \Big).
\]

\subsection*{Approach 1: Recursion}

\begin{lstlisting}
#include <bits/stdc++.h>
using namespace std;

int mcmRecursive(int i, int j, vector<int>& arr) {
    if (i == j) return 0;

    int minCost = INT_MAX;
    for (int k = i; k < j; k++) {
        int cost = mcmRecursive(i, k, arr) + mcmRecursive(k + 1, j, arr) +
                   arr[i - 1] * arr[k] * arr[j];
        minCost = min(minCost, cost);
    }
    return minCost;
}

int main() {
    vector<int> arr = {40, 20, 30, 10, 30};
    int n = arr.size() - 1;
    cout << mcmRecursive(1, n, arr) << endl; // 26000
    return 0;
}
\end{lstlisting}

\begin{complexitybox}[Approach 1 Complexity]
\textbf{Time.} Exponential -- roughly $O(2^n \cdot n)$ from the
Catalan-number-sized search space of parenthesizations.
\textbf{Space.} $O(n)$ recursion stack.
\end{complexitybox}

\subsection*{Approach 2: Memoization}

\begin{lstlisting}
#include <bits/stdc++.h>
using namespace std;

int mcmMemo(int i, int j, vector<int>& arr, vector<vector<int>>& memo) {
    if (i == j) return 0;
    if (memo[i][j] != -1) return memo[i][j];

    int minCost = INT_MAX;
    for (int k = i; k < j; k++) {
        int cost = mcmMemo(i, k, arr, memo) + mcmMemo(k + 1, j, arr, memo) +
                   arr[i - 1] * arr[k] * arr[j];
        minCost = min(minCost, cost);
    }
    return memo[i][j] = minCost;
}

int main() {
    vector<int> arr = {40, 20, 30, 10, 30};
    int n = arr.size() - 1;
    vector<vector<int>> memo(n + 1, vector<int>(n + 1, -1));
    cout << mcmMemo(1, n, arr, memo) << endl; // 26000
    return 0;
}
\end{lstlisting}

\begin{complexitybox}[Approach 2 Complexity]
\textbf{Time.} $O(n^2)$ distinct $(i,j)$ states, each doing an $O(n)$
scan over $k$: $O(n^3)$.
\textbf{Space.} $O(n^2)$ memo $+$ $O(n)$ recursion stack.
\end{complexitybox}

\subsection*{Approach 3: Tabulation}

\begin{ideabox}[Filling Order: by Increasing Range Length, Not by Index]
$f(i,j)$ depends only on \emph{strictly smaller} ranges $(i,k)$ and
$(k{+}1,j)$, both contained within $[i,j]$. So the table must be filled
in order of increasing range length ($j-i$), not in simple increasing
$i$ or $j$ order -- the first new fill-order strategy in this book.
\end{ideabox}

\begin{lstlisting}
#include <bits/stdc++.h>
using namespace std;

int mcmTabulation(vector<int>& arr) {
    int n = arr.size() - 1;
    vector<vector<int>> dp(n + 1, vector<int>(n + 1, 0));

    for (int len = 1; len < n; len++) {
        for (int i = 1; i + len <= n; i++) {
            int j = i + len;
            dp[i][j] = INT_MAX;
            for (int k = i; k < j; k++) {
                int cost = dp[i][k] + dp[k + 1][j] +
                           arr[i - 1] * arr[k] * arr[j];
                dp[i][j] = min(dp[i][j], cost);
            }
        }
    }
    return dp[1][n];
}

int main() {
    vector<int> arr = {40, 20, 30, 10, 30};
    cout << mcmTabulation(arr) << endl; // 26000
    return 0;
}
\end{lstlisting}

\subsection*{Dry run}

\dryrun{Take $\textit{arr} = [40,20,30,10,30]$ ($4$ matrices: $40{\times}20$,
$20{\times}30$, $30{\times}10$, $10{\times}30$).}

For the full range $(1,4)$, the best split is at $k=3$: multiply
matrices $1..3$ combined ($f(1,3)=14000$, itself optimally split at
$k=1$ inside that sub-range) with matrix $4$ alone ($f(4,4)=0$), joined
at cost $\textit{arr}[0]\cdot\textit{arr}[3]\cdot\textit{arr}[4] =
40\cdot10\cdot30=12000$. Total:
$f(1,4)=14000+0+12000=26000$, matching the expected answer (the
alternative splits at $k=1$ and $k=2$ both give strictly larger
totals, $36000$ and $69000$ respectively).

\begin{complexitybox}[Approach 3 Complexity]
\textbf{Time.} $O(n^3)$ -- $O(n^2)$ ranges, each scanning $O(n)$ split
points.
\textbf{Space.} $O(n^2)$ for the table. There is no meaningful
space-optimized version: every cell $(i,j)$ genuinely depends on
arbitrarily distant smaller ranges, not just a fixed nearby row, so the
full 2D table must be retained.
\end{complexitybox}

\begin{takeawaybox}
Interval DP's defining feature -- filling by range length, with an
inner search over every split point -- appears here for the first time
in this book, and every other section in this chapter reuses this exact
fill order with only the cost formula changed.
\end{takeawaybox}
